\documentclass[aspectratio=169, 12pt]{beamer}
\usepackage[utf8]{inputenc}
\usepackage[T1]{fontenc}
\usepackage[brazilian]{babel}
\usepackage{graphicx}
\usepackage{amsmath}
\usepackage{booktabs}
\usepackage{siunitx}
\usepackage{tikz}
\usetikzlibrary{arrows.meta, positioning, shapes.geometric, calc, decorations.pathreplacing}
\usepackage{fontawesome5}
\usepackage{xcolor}
\usepackage{multicol}
\usepackage{caption}

% =============================================
% TEMA E CORES
% =============================================
\usetheme{Madrid}
\definecolor{uerjazul}{RGB}{0, 51, 102}
\definecolor{uerjvermelho}{RGB}{153, 0, 0}
\definecolor{cinzaclaro}{RGB}{240, 240, 240}
\definecolor{verdepositivo}{RGB}{34, 139, 34}
\definecolor{laranjaalerta}{RGB}{255, 140, 0}

\setbeamercolor{structure}{fg=uerjazul}
\setbeamercolor{frametitle}{bg=uerjazul, fg=white}
\setbeamercolor{title}{bg=uerjazul, fg=white}
\setbeamercolor{block title}{bg=uerjazul, fg=white}
\setbeamercolor{block body}{bg=cinzaclaro}
\setbeamercolor{block title alerted}{bg=uerjvermelho, fg=white}
\setbeamercolor{block title example}{bg=verdepositivo, fg=white}

% =============================================
% INFORMAÇÕES DO DOCUMENTO
% =============================================
\title[IA na Detecção de Corrosão]{%
    Inteligência Artificial Aplicada à\\
    Detecção de Corrosão em Estruturas de Aço}
\subtitle{Como Ensinamos um Computador a ``Enxergar'' Ferrugem}

\author[Heitor O. Gonçalves]{%
    \textbf{Heitor Oliveira Gonçalves}\\
    \small{Universidade Católica de Petrópolis (UCP)}\\
    \small{Orientadores: Prof. Darlan Porto $\cdot$ Prof. Renato Amaral}
}

\institute[UCP/UERJ]{%
    Apresentação na\\
    \textbf{Universidade do Estado do Rio de Janeiro (UERJ)}\\
    Faculdade de Engenharia
}

\date{2026}

\begin{document}

% =============================================
% SLIDE 1: CAPA
% =============================================
\begin{frame}[plain]
    \titlepage
\end{frame}

% =============================================
% SLIDE 2: ROTEIRO
% =============================================
\begin{frame}{O que vamos conversar hoje}
    \begin{columns}[T]
        \column{0.02\textwidth}
        \column{0.96\textwidth}
        \begin{enumerate}
            \item \textbf{O Problema} --- Por que a corrosão é um desafio?
            \item \textbf{A Ideia} --- Como a IA pode ajudar?
            \item \textbf{Como Funciona} --- O que é Deep Learning na prática?
            \item \textbf{O Que Construímos} --- Ferramentas e arquitetura do sistema
            \item \textbf{Comparação de Técnicas} --- U-Net vs Attention U-Net vs Classificadores
            \item \textbf{Resultados} --- O computador realmente aprendeu?
            \item \textbf{Na Prática} --- O que isso muda para o engenheiro civil?
        \end{enumerate}
    \end{columns}
\end{frame}

% =============================================
% PARTE 1: O PROBLEMA
% =============================================
\section{O Problema}

\begin{frame}{O Inimigo Silencioso da Engenharia Civil}
    \begin{columns}[T]
        \column{0.55\textwidth}
        \begin{block}{A Corrosão em Números}
            \begin{itemize}
                \item Perdas globais: \textbf{US\$ 2,5 trilhões/ano}
                \item Equivale a \textbf{3,4\% do PIB mundial}
                \item No Brasil: bilhões em manutenção e substituição
                \item Causa principal de interdição de pontes e viadutos
            \end{itemize}
        \end{block}

        \vspace{0.3cm}
        \begin{alertblock}{A Pergunta Central}
            \textit{E se pudéssemos detectar a corrosão de forma\\automática, objetiva e antes que ela se torne crítica?}
        \end{alertblock}

        \column{0.42\textwidth}
        \begin{block}{Como é feito hoje?}
            \begin{itemize}
                \item[\faEye] Inspeção \textbf{visual} manual
                \item[\faUserClock] Depende da \textbf{experiência} do inspetor
                \item[\faClock] Processo \textbf{lento} e caro
                \item[\faExclamationTriangle] Resultado \textbf{subjetivo}
                \item[\faMapMarkerAlt] Áreas de \textbf{difícil acesso}
            \end{itemize}
        \end{block}
    \end{columns}
\end{frame}

\begin{frame}{O Desafio: Dois Inspetores, Duas Opiniões}
    \begin{block}{O Problema da Subjetividade}
        Imagine dois engenheiros inspecionando a mesma viga de aço:
    \end{block}

    \vspace{0.3cm}

    \begin{columns}[T]
        \column{0.45\textwidth}
        \begin{exampleblock}{Inspetor A}
            \begin{itemize}
                \item ``Corrosão \textbf{moderada}''
                \item ``Precisa de acompanhamento''
                \item ``Não é urgente''
            \end{itemize}
        \end{exampleblock}

        \column{0.45\textwidth}
        \begin{alertblock}{Inspetor B}
            \begin{itemize}
                \item ``Corrosão \textbf{severa}''
                \item ``Intervenção imediata''
                \item ``Risco estrutural''
            \end{itemize}
        \end{alertblock}
    \end{columns}

    \vspace{0.5cm}

    \begin{block}{A Solução}
        Um sistema de IA não tem ``dia ruim''. Ele analisa \textbf{cada pixel} da imagem com o \textbf{mesmo critério}, sempre. Sem cansaço, sem pressa, sem subjetividade.
    \end{block}
\end{frame}

% =============================================
% PARTE 2: A IDEIA
% =============================================
\section{A Ideia --- Como a IA Pode Ajudar?}

\begin{frame}{De Olho Humano para Olho Digital}
    \begin{block}{A Analogia}
        Pense em como \textbf{você} aprendeu a reconhecer corrosão:
    \end{block}

    \vspace{0.2cm}

    \begin{columns}[T]
        \column{0.48\textwidth}
        \begin{block}{\faGraduationCap\hspace{0.2cm} Aprendizado Humano}
            \begin{enumerate}
                \item Viu \textbf{exemplos} na faculdade
                \item Professor \textbf{apontou} onde estava a corrosão
                \item Com \textbf{prática}, aprendeu a identificar
                \item Hoje reconhece padrões \textbf{automaticamente}
            \end{enumerate}
        \end{block}

        \column{0.48\textwidth}
        \begin{block}{\faRobot\hspace{0.2cm} Aprendizado da Máquina}
            \begin{enumerate}
                \item Mostramos \textbf{centenas de fotos}
                \item \textbf{Marcamos} onde está a corrosão
                \item A rede neural \textbf{ajusta} seus pesos
                \item O modelo reconhece padrões \textbf{sozinho}
            \end{enumerate}
        \end{block}
    \end{columns}

    \vspace{0.4cm}

    \begin{alertblock}{Diferença Fundamental}
        O engenheiro vê dezenas de estruturas por semana.\\
        O computador ``vê'' \textbf{800 imagens em minutos} e aprende com \textbf{cada pixel}.
    \end{alertblock}
\end{frame}

\begin{frame}{O Que é Deep Learning? (Sem Jargão)}
    \begin{block}{Explicação Simples}
        Deep Learning é como um \textbf{filtro inteligente de camadas}:
    \end{block}

    \vspace{0.3cm}

    \begin{center}
    \begin{tikzpicture}[
        node distance=0.6cm,
        box/.style={rectangle, rounded corners, draw=uerjazul, fill=uerjazul!10,
                     minimum width=3.8cm, minimum height=0.9cm, align=center, font=\small},
        arrow/.style={-{Stealth[length=3mm]}, thick, uerjazul}
    ]
        \node[box] (input) {\textbf{Camada 1:} Detecta bordas\\e cores básicas};
        \node[box, below=of input] (mid1) {\textbf{Camada 2:} Reconhece\\texturas e formas};
        \node[box, below=of mid1] (mid2) {\textbf{Camada 3:} Identifica\\padrões de corrosão};
        \node[box, below=of mid2, fill=verdepositivo!15, draw=verdepositivo] (output) {\textbf{Camada Final:} ``Aqui tem\\corrosão!'' (ou não)};

        \draw[arrow] (input) -- (mid1);
        \draw[arrow] (mid1) -- (mid2);
        \draw[arrow] (mid2) -- (output);

        \node[left=1.5cm of input, font=\small, align=center] (foto) {\faCamera\\Foto da\\estrutura};
        \node[right=1.5cm of output, font=\small, align=center] (result) {\faCheckCircle\\Mapa de\\corrosão};

        \draw[arrow, dashed] (foto) -- (input);
        \draw[arrow, dashed] (output) -- (result);
    \end{tikzpicture}
    \end{center}

    \vspace{0.2cm}

    \begin{block}{}
        Cada camada aprende algo mais complexo. É como olhar com uma \textbf{lupa cada vez mais potente}.
    \end{block}
\end{frame}

% =============================================
% PARTE 3: COMO FUNCIONA
% =============================================
\section{Como Funciona na Prática}

\begin{frame}{Passo a Passo: Do Zero ao Sistema Funcionando}
    \begin{center}
    \begin{tikzpicture}[
        node distance=0.15cm and 0.6cm,
        step/.style={rectangle, rounded corners=4pt, draw=uerjazul, fill=uerjazul!8,
                     minimum width=2.4cm, minimum height=1.5cm, align=center, font=\scriptsize},
        arrow/.style={-{Stealth[length=2.5mm]}, thick, uerjazul!70}
    ]
        \node[step] (s1) {\faCamera\\\textbf{1. Coleta}\\Fotografar\\as vigas};
        \node[step, right=of s1] (s2) {\faPaintBrush\\\textbf{2. Anotação}\\Marcar onde\\tem corrosão};
        \node[step, right=of s2] (s3) {\faCogs\\\textbf{3. Preparo}\\Redimensionar\\e organizar};
        \node[step, right=of s3] (s4) {\faBrain\\\textbf{4. Treino}\\Ensinar a\\rede neural};
        \node[step, right=of s4, fill=verdepositivo!10, draw=verdepositivo] (s5) {\faCheckDouble\\\textbf{5. Teste}\\Avaliar com\\imagens novas};

        \draw[arrow] (s1) -- (s2);
        \draw[arrow] (s2) -- (s3);
        \draw[arrow] (s3) -- (s4);
        \draw[arrow] (s4) -- (s5);
    \end{tikzpicture}
    \end{center}

    \vspace{0.5cm}

    \begin{columns}[T]
        \column{0.48\textwidth}
        \begin{block}{O que é ``ensinar'' a rede?}
            A rede vê a foto e tenta adivinhar onde está a corrosão. Comparamos com a resposta correta (a marcação do especialista) e ajustamos os ``pesos'' internos. Repetimos isso \textbf{milhares de vezes}.
        \end{block}

        \column{0.48\textwidth}
        \begin{block}{O que é ``testar''?}
            Mostramos fotos que a rede \textbf{nunca viu antes}. Se ela acerta, significa que realmente \textit{aprendeu} --- e não apenas \textit{decorou}.
        \end{block}
    \end{columns}
\end{frame}

\begin{frame}{Etapa 1: Construindo o Dataset}
    \begin{block}{As ``Fichas de Estudo'' da Inteligência Artificial}
        Assim como um aluno precisa de livros, a IA precisa de \textbf{dados}.
    \end{block}

    \vspace{0.3cm}

    \begin{columns}[T]
        \column{0.48\textwidth}
        \begin{block}{Nosso Dataset}
            \begin{itemize}
                \item \textbf{800 imagens} de vigas W de aço
                \item Aço ASTM A572 Grau 50
                \item Resolução: 256 $\times$ 256 pixels
                \item Imagens em escala de cinza
                \item Diferentes níveis de corrosão
            \end{itemize}
        \end{block}

        \column{0.48\textwidth}
        \begin{block}{Divisão Estratégica}
            \begin{itemize}
                \item[\textcolor{uerjazul}{\textbf{70\%}}] para \textbf{treino} (560 imagens)\\
                    {\scriptsize A IA estuda com estas}
                \item[\textcolor{laranjaalerta}{\textbf{15\%}}] para \textbf{validação} (120 imagens)\\
                    {\scriptsize ``Simulado'' durante o estudo}
                \item[\textcolor{verdepositivo}{\textbf{15\%}}] para \textbf{teste} (120 imagens)\\
                    {\scriptsize ``Prova final'' --- nunca vista}
            \end{itemize}
        \end{block}
    \end{columns}

    \vspace{0.3cm}

    \begin{alertblock}{Por que dividir?}
        Se a IA estudasse e fizesse prova com as mesmas imagens, seria como colar na prova. A divisão garante que ela realmente \textbf{generalizou} o conhecimento.
    \end{alertblock}
\end{frame}

\begin{frame}{Etapa 2: As ``Gabaritos'' --- Máscaras de Segmentação}
    \begin{columns}[T]
        \column{0.55\textwidth}
        \begin{block}{O que é uma máscara?}
            Para cada foto de uma viga, um \textbf{especialista em patologia estrutural} pintou manualmente as regiões corroídas, pixel por pixel.
        \end{block}

        \vspace{0.2cm}

        \begin{block}{Como funciona}
            \begin{itemize}
                \item \textbf{Pixel branco} = corrosão presente
                \item \textbf{Pixel preto} = aço saudável
                \item Ferramenta: CVAT (Computer Vision Annotation Tool)
                \item Concordância entre anotadores: \textbf{87\%} (Kappa de Fleiss)
            \end{itemize}
        \end{block}

        \column{0.42\textwidth}
        \begin{block}{Visualização Conceitual}
            \begin{center}
            \begin{tikzpicture}[scale=0.85]
                % Imagem original
                \draw[thick, uerjazul] (0,0) rectangle (3,2.2);
                \node[font=\scriptsize, above] at (1.5,2.2) {\textbf{Foto Original}};
                \fill[uerjazul!20] (0,0) rectangle (3,2.2);
                % Manchas de corrosão simuladas
                \fill[uerjvermelho!50] (0.5,0.8) circle (0.4);
                \fill[uerjvermelho!40] (2.0,1.4) ellipse (0.5 and 0.3);
                \fill[uerjvermelho!45] (1.2,0.3) circle (0.25);

                % Seta
                \draw[-{Stealth}, thick] (3.3,1.1) -- (3.9,1.1);

                % Máscara
                \draw[thick, uerjazul] (4.2,0) rectangle (7.2,2.2);
                \node[font=\scriptsize, above] at (5.7,2.2) {\textbf{Máscara (Gabarito)}};
                \fill[black] (4.2,0) rectangle (7.2,2.2);
                % Regiões marcadas
                \fill[white] (4.7,0.8) circle (0.4);
                \fill[white] (6.2,1.4) ellipse (0.5 and 0.3);
                \fill[white] (5.4,0.3) circle (0.25);
            \end{tikzpicture}
            \end{center}
        \end{block}

        \vspace{0.2cm}
        {\scriptsize \textit{As regiões em branco na máscara são exatamente onde o especialista identificou corrosão.}}
    \end{columns}
\end{frame}

\begin{frame}{Etapa 3: Data Augmentation --- ``Multiplicando'' os Dados}
    \begin{block}{O Problema}
        800 imagens é pouco para Deep Learning. Solução: \textbf{criar variações} de cada imagem.
    \end{block}

    \vspace{0.3cm}

    \begin{center}
    \begin{tikzpicture}[
        img/.style={rectangle, draw=uerjazul, minimum width=1.6cm, minimum height=1.2cm, align=center, font=\scriptsize, fill=uerjazul!5},
        arrow/.style={-{Stealth[length=2mm]}, uerjazul, thick}
    ]
        \node[img, minimum width=2cm, minimum height=1.6cm, fill=uerjazul!15] (orig) {\faImage\\Original};

        \node[img, right=1.2cm of orig, yshift=1.2cm] (r1) {\faSyncAlt\\Rotação\\$90^{\circ}$};
        \node[img, right=1.2cm of orig, yshift=0cm] (r2) {\faArrowsAltH\\Espelhamento\\horizontal};
        \node[img, right=1.2cm of orig, yshift=-1.2cm] (r3) {\faSun\\Variação de\\brilho};

        \draw[arrow] (orig.east) -- (r1.west);
        \draw[arrow] (orig.east) -- (r2.west);
        \draw[arrow] (orig.east) -- (r3.west);

        \node[right=0.3cm of r2, font=\small, align=left] {$\Rightarrow$ De 800\\passamos para\\milhares de\\exemplos};
    \end{tikzpicture}
    \end{center}

    \vspace{0.3cm}

    \begin{exampleblock}{Analogia para o engenheiro}
        É como estudar uma planta baixa e depois olhá-la de cabeça para baixo, espelhada, com iluminação diferente. O conteúdo é o mesmo --- mas o cérebro (e a IA) aprende a reconhecer de \textbf{qualquer ângulo}.
    \end{exampleblock}
\end{frame}

% =============================================
% PARTE 4: O QUE CONSTRUÍMOS
% =============================================
\section{O Que Construímos}

\begin{frame}{Ambiente de Desenvolvimento}
    \begin{columns}[T]
        \column{0.48\textwidth}
        \begin{block}{\faCogs\hspace{0.2cm} Software}
            \begin{itemize}
                \item \textbf{MATLAB} R2022a+
                \item Deep Learning Toolbox
                \item Image Processing Toolbox
                \item Computer Vision Toolbox
                \item Parallel Computing Toolbox (GPU)
            \end{itemize}
        \end{block}

        \vspace{0.2cm}

        \begin{block}{\faCode\hspace{0.2cm} Código Desenvolvido}
            \begin{itemize}
                \item \textbf{+7.000 linhas} de código MATLAB
                \item \textbf{50+ testes} automatizados
                \item \textbf{7 módulos} principais
                \item Arquitetura orientada a objetos
            \end{itemize}
        \end{block}

        \column{0.48\textwidth}
        \begin{block}{\faServer\hspace{0.2cm} Hardware}
            \begin{itemize}
                \item GPU NVIDIA com 4--24 GB VRAM
                \item 16 GB+ de RAM
                \item SSD para leitura rápida
            \end{itemize}
        \end{block}

        \vspace{0.2cm}

        \begin{block}{\faProjectDiagram\hspace{0.2cm} Arquitetura do Sistema}
            O sistema tem dois braços:
            \begin{enumerate}
                \item \textbf{Segmentação:} marca \textit{onde} está a corrosão (pixel a pixel)
                \item \textbf{Classificação:} diz \textit{quão grave} é (leve, moderada, severa)
            \end{enumerate}
        \end{block}
    \end{columns}
\end{frame}

\begin{frame}{Visão Geral do Pipeline}
    \begin{center}
    \begin{tikzpicture}[
        node distance=0.4cm and 0.2cm,
        box/.style={rectangle, rounded corners=3pt, draw=uerjazul, fill=uerjazul!8,
                     minimum width=2.5cm, minimum height=1cm, align=center, font=\scriptsize},
        boxg/.style={rectangle, rounded corners=3pt, draw=verdepositivo, fill=verdepositivo!8,
                     minimum width=2.5cm, minimum height=1cm, align=center, font=\scriptsize},
        boxr/.style={rectangle, rounded corners=3pt, draw=uerjvermelho, fill=uerjvermelho!8,
                     minimum width=2.5cm, minimum height=1cm, align=center, font=\scriptsize},
        arrow/.style={-{Stealth[length=2.5mm]}, thick, uerjazul!70}
    ]
        % Input
        \node[box, minimum width=2cm] (input) {\faCamera\\\textbf{Imagem}\\da viga};

        % Classificação
        \node[boxg, right=1.2cm of input, yshift=1cm] (class) {\faTachometerAlt\\\textbf{Classificação}\\ResNet50 / EfficientNet};
        \node[right=0.5cm of class, font=\scriptsize, align=left] (classout) {Nenhuma\\Leve\\Moderada\\Severa};

        % Segmentação
        \node[boxr, right=1.2cm of input, yshift=-1cm] (seg) {\faSearchPlus\\\textbf{Segmentação}\\U-Net / Att. U-Net};
        \node[right=0.5cm of seg, font=\scriptsize, align=left] (segout) {Mapa pixel\\a pixel de\\corrosão};

        % Relatório
        \node[box, right=2.5cm of input, yshift=-2.8cm, minimum width=4cm, fill=laranjaalerta!10, draw=laranjaalerta] (report) {\faFileAlt\hspace{0.2cm}\textbf{Relatório Completo:} localização + severidade};

        \draw[arrow] (input.east) -- ++(0.4,0) |- (class.west);
        \draw[arrow] (input.east) -- ++(0.4,0) |- (seg.west);
        \draw[arrow] (class.east) -- (classout.west);
        \draw[arrow] (seg.east) -- (segout.west);
        \draw[arrow, dashed, laranjaalerta] (classout.south) -- ++(0,-0.6) -| (report.north);
        \draw[arrow, dashed, laranjaalerta] (segout.south) -- ++(0,-0.2) -| (report.north);

        % Labels
        \node[above=0.15cm of class, font=\scriptsize\itshape, uerjazul] {Triagem rápida ($\sim$30ms)};
        \node[below=0.15cm of seg, font=\scriptsize\itshape, uerjazul] {Análise detalhada ($\sim$85--128ms)};
    \end{tikzpicture}
    \end{center}
\end{frame}

% =============================================
% PARTE 5: COMPARAÇÃO DE TÉCNICAS
% =============================================
\section{Comparação de Técnicas}

\begin{frame}{Duas Tarefas, Duas Abordagens}
    \begin{columns}[T]
        \column{0.48\textwidth}
        \begin{block}{\faSearchPlus\hspace{0.2cm} Segmentação Semântica}
            \textbf{``Onde \textit{exatamente} está a corrosão?''}
            \begin{itemize}
                \item Analisa \textbf{cada pixel}
                \item Gera um \textbf{mapa} (máscara)
                \item Como pintar com marcador
                \item Modelos: \textbf{U-Net} e \textbf{Attention U-Net}
            \end{itemize}

            \vspace{0.2cm}
            {\footnotesize \textit{Ideal para: engenheiro que precisa saber a extensão exata da corrosão}}
        \end{block}

        \column{0.48\textwidth}
        \begin{block}{\faTachometerAlt\hspace{0.2cm} Classificação}
            \textbf{``Qual o \textit{grau} de corrosão?''}
            \begin{itemize}
                \item Analisa a \textbf{imagem toda}
                \item Dá uma \textbf{nota} (classe)
                \item Como dar uma avaliação geral
                \item Modelos: \textbf{ResNet50}, \textbf{EfficientNet}, \textbf{CNN}
            \end{itemize}

            \vspace{0.2cm}
            {\footnotesize \textit{Ideal para: triagem rápida de centenas de fotos de inspeção}}
        \end{block}
    \end{columns}

    \vspace{0.5cm}

    \begin{exampleblock}{Analogia Médica}
        Segmentação = \textbf{ressonância magnética} (mostra exatamente onde está o problema).\\
        Classificação = \textbf{exame de sangue} (diz se está normal, alterado ou crítico).
    \end{exampleblock}
\end{frame}

\begin{frame}{U-Net: A Arquitetura Base}
    \begin{columns}[T]
        \column{0.55\textwidth}
        \begin{block}{Como funciona a U-Net?}
            Criada em 2015 para imagens médicas, funciona em duas etapas:

            \vspace{0.2cm}

            \textbf{1. Encoder (lado esquerdo do U):}
            \begin{itemize}
                \item ``Comprime'' a imagem
                \item Extrai características cada vez mais abstratas
                \item De 256$\times$256 $\rightarrow$ 16$\times$16
            \end{itemize}

            \vspace{0.2cm}

            \textbf{2. Decoder (lado direito do U):}
            \begin{itemize}
                \item ``Reconstrói'' a imagem
                \item Marca pixel por pixel
                \item De 16$\times$16 $\rightarrow$ 256$\times$256
            \end{itemize}
        \end{block}

        \column{0.42\textwidth}
        \begin{block}{Formato em ``U''}
            \begin{center}
            \begin{tikzpicture}[scale=0.55, transform shape,
                block/.style={rectangle, draw=uerjazul, fill=uerjazul!15, minimum width=0.5cm}
            ]
                % Encoder
                \node[block, minimum height=2.5cm] (e1) at (0,0) {};
                \node[block, minimum height=2cm] (e2) at (1,-0.25) {};
                \node[block, minimum height=1.5cm] (e3) at (2,-0.5) {};
                \node[block, minimum height=1cm, fill=uerjvermelho!20, draw=uerjvermelho] (bottom) at (3,-0.75) {};
                % Decoder
                \node[block, minimum height=1.5cm] (d3) at (4,-0.5) {};
                \node[block, minimum height=2cm] (d2) at (5,-0.25) {};
                \node[block, minimum height=2.5cm, fill=verdepositivo!15, draw=verdepositivo] (d1) at (6,0) {};

                % Arrows down
                \draw[-{Stealth}, thick] (e1.east) -- (e2.west);
                \draw[-{Stealth}, thick] (e2.east) -- (e3.west);
                \draw[-{Stealth}, thick] (e3.east) -- (bottom.west);
                % Arrows up
                \draw[-{Stealth}, thick] (bottom.east) -- (d3.west);
                \draw[-{Stealth}, thick] (d3.east) -- (d2.west);
                \draw[-{Stealth}, thick] (d2.east) -- (d1.west);

                % Skip connections
                \draw[-{Stealth}, dashed, uerjazul!50, thick] (e1.north) .. controls (3,2.5) .. (d1.north);
                \draw[-{Stealth}, dashed, uerjazul!50, thick] (e2.north) .. controls (3,1.8) .. (d2.north);
                \draw[-{Stealth}, dashed, uerjazul!50, thick] (e3.north) .. controls (3,1.2) .. (d3.north);

                % Labels
                \node[below=0.3cm, font=\tiny] at (0,-1.25) {Entrada};
                \node[below=0.3cm, font=\tiny] at (3,-1.25) {Compressão};
                \node[below=0.3cm, font=\tiny] at (6,-1.25) {Saída};
                \node[above=0.1cm, font=\tiny\itshape, uerjazul!70] at (3,2.7) {skip connections};
            \end{tikzpicture}
            \end{center}
        \end{block}

        \vspace{0.2cm}

        \begin{block}{Skip Connections}
            {\scriptsize As \textbf{setas tracejadas} levam detalhes finos do encoder direto para o decoder --- como atalhos que preservam informações importantes.}
        \end{block}
    \end{columns}
\end{frame}

\begin{frame}{Attention U-Net: A Evolução}
    \begin{block}{O Que Muda?}
        A Attention U-Net adiciona ``portões de atenção'' (\textit{attention gates}) que ensinam o modelo a \textbf{prestar mais atenção} nas regiões importantes e \textbf{ignorar} as irrelevantes.
    \end{block}

    \vspace{0.3cm}

    \begin{columns}[T]
        \column{0.48\textwidth}
        \begin{exampleblock}{Analogia para o Engenheiro}
            Imagine que você está inspecionando uma viga de 10 metros:
            \begin{itemize}
                \item A U-Net olha \textbf{tudo igualmente}
                \item A Attention U-Net \textbf{foca automaticamente} nas regiões suspeitas, como um inspetor experiente que já sabe onde procurar
            \end{itemize}
        \end{exampleblock}

        \column{0.48\textwidth}
        \begin{block}{Dados Comparativos}
            \begin{tabular}{lcc}
                \toprule
                & \textbf{U-Net} & \textbf{Att. U-Net} \\
                \midrule
                Camadas & 70 & 94 \\
                Parâmetros & 23,5M & 31,0M \\
                Épocas & 50 & 60 \\
                Memória GPU & 2,2 GB & 2,8 GB \\
                \bottomrule
            \end{tabular}

            \vspace{0.3cm}

            {\scriptsize \textit{Mais parâmetros = mais capacidade de aprender padrões complexos, mas exige mais poder computacional.}}
        \end{block}
    \end{columns}
\end{frame}

\begin{frame}{Classificação: ResNet50, EfficientNet e CNN Customizada}
    \begin{block}{A Outra Metade do Sistema}
        Enquanto a segmentação mostra \textit{onde}, a classificação diz \textit{quão grave}:
    \end{block}

    \vspace{0.2cm}

    \begin{columns}[T]
        \column{0.32\textwidth}
        \begin{block}{\footnotesize ResNet50}
            \begin{itemize}
                \item[\faAward] \textbf{Mais preciso}
                \item 25 milhões de parâmetros
                \item Transfer learning do ImageNet
                \item Acurácia: \textbf{94,2\%}
                \item $\sim$45ms por imagem
            \end{itemize}
        \end{block}

        \column{0.32\textwidth}
        \begin{block}{\footnotesize EfficientNet-B0}
            \begin{itemize}
                \item[\faBalanceScale] \textbf{Melhor equilíbrio}
                \item 5 milhões de parâmetros
                \item Arquitetura otimizada
                \item Acurácia: \textbf{92,8\%}
                \item $\sim$30ms por imagem
            \end{itemize}
        \end{block}

        \column{0.32\textwidth}
        \begin{block}{\footnotesize CNN Customizada}
            \begin{itemize}
                \item[\faBolt] \textbf{Mais rápido}
                \item 2 milhões de parâmetros
                \item Feito sob medida
                \item Acurácia: \textbf{89,5\%}
                \item $\sim$20ms por imagem
            \end{itemize}
        \end{block}
    \end{columns}

    \vspace{0.3cm}

    \begin{alertblock}{O que é Transfer Learning?}
        A ResNet50 e a EfficientNet já ``aprenderam'' a reconhecer formas e texturas com 14 milhões de fotos (ImageNet). Nós apenas \textbf{refinamos} esse conhecimento para corrosão. É como contratar um inspetor experiente e só treiná-lo no tipo específico de estrutura.
    \end{alertblock}
\end{frame}

\begin{frame}{As 4 Classes de Severidade}
    \begin{center}
    \begin{tikzpicture}[
        classe/.style={rectangle, rounded corners=5pt, minimum width=2.5cm, minimum height=1.8cm,
                       align=center, font=\small, draw, thick}
    ]
        \node[classe, fill=verdepositivo!15, draw=verdepositivo] (c0) at (0,0) {
            \textbf{Nenhuma}\\
            {\scriptsize 0\% corrosão}\\
            {\scriptsize \faCheckCircle\hspace{0.1cm} Saudável}
        };
        \node[classe, fill=yellow!15, draw=yellow!80!black] (c1) at (3.5,0) {
            \textbf{Leve}\\
            {\scriptsize $<$ 10\%}\\
            {\scriptsize \faEye\hspace{0.1cm} Monitorar}
        };
        \node[classe, fill=laranjaalerta!15, draw=laranjaalerta] (c2) at (7,0) {
            \textbf{Moderada}\\
            {\scriptsize 10--30\%}\\
            {\scriptsize \faTools\hspace{0.1cm} Planejar}
        };
        \node[classe, fill=uerjvermelho!15, draw=uerjvermelho] (c3) at (10.5,0) {
            \textbf{Severa}\\
            {\scriptsize $>$ 30\%}\\
            {\scriptsize \faExclamationTriangle\hspace{0.1cm} Urgente}
        };

        % Seta de severidade
        \draw[-{Stealth}, very thick, uerjvermelho!50] (0,-1.4) -- (10.5,-1.4)
            node[midway, below, font=\small] {Severidade crescente};
    \end{tikzpicture}
    \end{center}

    \vspace{0.3cm}

    \begin{block}{Como o sistema decide a classe?}
        \begin{enumerate}
            \item A máscara de segmentação conta os \textbf{pixels brancos} (corrosão)
            \item Calcula a \textbf{porcentagem} da área corroída
            \item Classifica automaticamente na faixa correspondente
            \item Gera o \textbf{rótulo} (label) usado para treinar os classificadores
        \end{enumerate}
    \end{block}
\end{frame}

% =============================================
% PARTE 6: RESULTADOS
% =============================================
\section{Resultados}

\begin{frame}{Segmentação: U-Net vs Attention U-Net}
    \begin{table}
        \centering
        \small
        \begin{tabular}{l c c c}
            \toprule
            \textbf{Métrica} & \textbf{U-Net} & \textbf{Att. U-Net} & \textbf{Melhoria} \\
            \midrule
            IoU (sobreposição) & 0,743 & \textbf{0,764} & \textcolor{verdepositivo}{+2,9\%} \\
            Dice (concordância) & 0,830 & \textbf{0,868} & \textcolor{verdepositivo}{+4,5\%} \\
            Acurácia por pixel & 0,882 & \textbf{0,918} & \textcolor{verdepositivo}{+4,1\%} \\
            Precisão & 0,811 & \textbf{0,856} & \textcolor{verdepositivo}{+5,6\%} \\
            Recall (sensibilidade) & 0,750 & \textbf{0,799} & \textcolor{verdepositivo}{+6,6\%} \\
            F1-Score & 0,774 & \textbf{0,816} & \textcolor{verdepositivo}{+5,5\%} \\
            \bottomrule
        \end{tabular}
    \end{table}

    \vspace{0.3cm}

    \begin{block}{O que esses números significam?}
        \begin{itemize}
            \item \textbf{IoU $\approx$ 0,76}: de cada 100 pixels que o modelo marca como corrosão, \textbf{76 estão corretos} (e vice-versa)
            \item \textbf{Acurácia 91,8\%}: a cada 1.000 pixels, o modelo acerta a classificação de \textbf{918}
            \item Todos os resultados com \textbf{significância estatística} (p $<$ 0,01)
        \end{itemize}
    \end{block}
\end{frame}

\begin{frame}{Classificação: Os Três Modelos}
    \begin{table}
        \centering
        \small
        \begin{tabular}{l c c c}
            \toprule
            \textbf{Modelo} & \textbf{Acurácia} & \textbf{F1-Score} & \textbf{Velocidade} \\
            \midrule
            \textbf{ResNet50} & \textcolor{verdepositivo}{\textbf{94,2\%}} & 0,940 & 45ms/imagem \\
            EfficientNet-B0 & 92,8\% & 0,930 & 30ms/imagem \\
            CNN Customizada & 89,5\% & 0,900 & \textcolor{verdepositivo}{\textbf{20ms/imagem}} \\
            \bottomrule
        \end{tabular}
    \end{table}

    \vspace{0.3cm}

    \begin{columns}[T]
        \column{0.48\textwidth}
        \begin{block}{Para o engenheiro, isso significa:}
            \begin{itemize}
                \item A cada \textbf{100 fotos} analisadas pela ResNet50, \textbf{94 são classificadas corretamente}
                \item Processamento de \textbf{1.000 fotos em menos de 1 minuto}
                \item Sem cansaço, sem subjetividade
            \end{itemize}
        \end{block}

        \column{0.48\textwidth}
        \begin{block}{Qual escolher?}
            \begin{itemize}
                \item[\faAward] \textbf{Máxima precisão:} ResNet50
                \item[\faBalanceScale] \textbf{Equilíbrio:} EfficientNet-B0
                \item[\faMobileAlt] \textbf{Dispositivo portátil:} CNN Customizada
            \end{itemize}
        \end{block}
    \end{columns}
\end{frame}

\begin{frame}{Velocidade: Quanto Tempo Leva?}
    \begin{block}{Tempo de Processamento}
        \begin{center}
        \begin{tikzpicture}[xscale=0.95]
            % Barras
            \fill[uerjazul!60] (0,0) rectangle (2,0.6);
            \node[right, font=\small] at (2.1,0.3) {CNN: \textbf{20ms}};
            \node[left, font=\small, white] at (1.9,0.3) {};

            \fill[uerjazul!70] (0,1) rectangle (3,1.6);
            \node[right, font=\small] at (3.1,1.3) {EfficientNet: \textbf{30ms}};

            \fill[uerjazul!80] (0,2) rectangle (4.5,2.6);
            \node[right, font=\small] at (4.6,2.3) {ResNet50: \textbf{45ms}};

            \fill[uerjvermelho!50] (0,3) rectangle (8.5,3.6);
            \node[right, font=\small] at (8.6,3.3) {U-Net: \textbf{84ms}};

            \fill[uerjvermelho!60] (0,4) rectangle (12.7,4.6);
            \node[right, font=\footnotesize] at (12.0,4.3) {\textbf{128ms}};
            \node[left, font=\small] at (11.9,4.3) {Attention U-Net:};

            % Eixo
            \node[font=\scriptsize] at (6,-0.4) {Tempo de inferência por imagem (milissegundos)};
        \end{tikzpicture}
        \end{center}
    \end{block}

    \vspace{0.2cm}

    \begin{exampleblock}{Contexto Prático}
        \begin{itemize}
            \item Um piscar de olhos leva $\sim$300ms --- todos os modelos são \textbf{mais rápidos que um piscar}
            \item Inspeção completa de 500 fotos: \textbf{menos de 2 minutos} (vs. dias para inspeção manual)
            \item Viável para uso em \textbf{tempo real} com drones
        \end{itemize}
    \end{exampleblock}
\end{frame}

% =============================================
% PARTE 7: NA PRÁTICA
% =============================================
\section{Impacto na Engenharia Civil}

\begin{frame}{O Que Muda Para o Engenheiro?}
    \begin{columns}[T]
        \column{0.48\textwidth}
        \begin{alertblock}{\faTimesCircle\hspace{0.2cm} Antes (Manual)}
            \begin{itemize}
                \item Inspetor escala a estrutura
                \item Avaliação visual subjetiva
                \item Relatório escrito à mão
                \item Dias para cobrir uma ponte
                \item Resultado depende do inspetor
                \item Custo elevado por inspeção
            \end{itemize}
        \end{alertblock}

        \column{0.48\textwidth}
        \begin{exampleblock}{\faCheckCircle\hspace{0.2cm} Depois (IA)}
            \begin{itemize}
                \item Drone captura imagens
                \item IA analisa automaticamente
                \item Relatório digital com mapa
                \item Minutos para centenas de fotos
                \item Resultado consistente e objetivo
                \item Custo reduzido por escala
            \end{itemize}
        \end{exampleblock}
    \end{columns}

    \vspace{0.5cm}

    \begin{block}{O Engenheiro \textbf{NÃO} é Substituído}
        A IA é uma \textbf{ferramenta} de apoio. O engenheiro continua tomando as decisões --- mas agora com dados objetivos, mapas precisos e classificações consistentes. É como trocar uma calculadora manual por uma planilha: o conhecimento continua sendo do profissional.
    \end{block}
\end{frame}

\begin{frame}{Aplicações Concretas}
    \begin{columns}[T]
        \column{0.48\textwidth}
        \begin{block}{\faBridgeWater\hspace{0.2cm} Pontes e Viadutos}
            \begin{itemize}
                \item Monitoramento periódico com drone
                \item Detecção precoce de degradação
                \item Priorização de manutenção
                \item Histórico digital da evolução
            \end{itemize}
        \end{block}

        \vspace{0.2cm}

        \begin{block}{\faBuilding\hspace{0.2cm} Edifícios Industriais}
            \begin{itemize}
                \item Galpões com estrutura metálica
                \item Áreas de difícil acesso
                \item Inspeções mais frequentes
                \item Custo reduzido
            \end{itemize}
        \end{block}

        \column{0.48\textwidth}
        \begin{block}{\faBroadcastTower\hspace{0.2cm} Torres e Infraestrutura}
            \begin{itemize}
                \item Torres de transmissão
                \item Postes metálicos
                \item Estruturas offshore
                \item Plantas industriais
            \end{itemize}
        \end{block}

        \vspace{0.2cm}

        \begin{block}{\faChartLine\hspace{0.2cm} Manutenção Preditiva}
            \begin{itemize}
                \item Acompanhamento temporal
                \item Previsão de vida útil
                \item Otimização de investimentos
                \item Conformidade normativa
            \end{itemize}
        \end{block}
    \end{columns}
\end{frame}

\begin{frame}{Fluxo de Trabalho Integrado com Drone}
    \begin{center}
    \begin{tikzpicture}[
        node distance=0.4cm and 0.2cm,
        step/.style={rectangle, rounded corners=4pt, draw=uerjazul, fill=uerjazul!8,
                     minimum width=2.2cm, minimum height=1.6cm, align=center, font=\scriptsize},
        arrow/.style={-{Stealth[length=2.5mm]}, thick, uerjazul!70}
    ]
        \node[step] (s1) {\faPlane\\\textbf{Drone}\\sobrevoa a\\estrutura};
        \node[step, right=of s1] (s2) {\faCamera\\\textbf{Captura}\\centenas de\\fotos em HD};
        \node[step, right=of s2] (s3) {\faUpload\\\textbf{Upload}\\envia ao\\servidor};
        \node[step, right=of s3, draw=verdepositivo, fill=verdepositivo!8] (s4) {\faBrain\\\textbf{IA Analisa}\\classifica e\\segmenta};
        \node[step, right=of s4, draw=laranjaalerta, fill=laranjaalerta!8] (s5) {\faFileAlt\\\textbf{Relatório}\\mapas e\\prioridades};

        \draw[arrow] (s1) -- (s2);
        \draw[arrow] (s2) -- (s3);
        \draw[arrow] (s3) -- (s4);
        \draw[arrow] (s4) -- (s5);
    \end{tikzpicture}
    \end{center}

    \vspace{0.5cm}

    \begin{columns}[T]
        \column{0.48\textwidth}
        \begin{block}{Benefícios Quantificados}
            \begin{itemize}
                \item \textbf{94,2\%} de acurácia na classificação
                \item \textbf{91,8\%} de acurácia na segmentação
                \item Processamento em \textbf{tempo real}
                \item Redução de \textbf{custos operacionais}
            \end{itemize}
        \end{block}

        \column{0.48\textwidth}
        \begin{block}{Segurança}
            \begin{itemize}
                \item Reduz exposição humana a riscos
                \item Acessa locais perigosos ou inacessíveis
                \item Documentação automática
                \item Rastreabilidade total
            \end{itemize}
        \end{block}
    \end{columns}
\end{frame}

% =============================================
% PARTE 8: MÉTRICAS EXPLICADAS
% =============================================
\section{Entendendo as Métricas}

\begin{frame}{Métricas de Avaliação --- Explicação Visual}
    \begin{block}{Como sabemos se o modelo é bom?}
        Usamos métricas que comparam a ``resposta'' da IA com o ``gabarito'' do especialista:
    \end{block}

    \vspace{0.3cm}

    \begin{columns}[T]
        \column{0.48\textwidth}
        \begin{block}{\textbf{IoU} (Intersection over Union)}
            {\small Quanto a resposta da IA ``sobrepõe'' o gabarito.}

            \begin{center}
            \begin{tikzpicture}[scale=0.7]
                \fill[uerjazul!20] (0,0) circle (1.2);
                \fill[uerjvermelho!20] (1.2,0) circle (1.2);
                \begin{scope}
                    \clip (0,0) circle (1.2);
                    \fill[verdepositivo!40] (1.2,0) circle (1.2);
                \end{scope}
                \node[font=\tiny] at (-0.4,0) {IA};
                \node[font=\tiny] at (1.6,0) {Real};
                \node[font=\tiny, verdepositivo] at (0.6,0) {IoU};
            \end{tikzpicture}
            \end{center}

            {\scriptsize \textbf{0,76} = sobreposição de 76\%. Excelente.}
        \end{block}

        \column{0.48\textwidth}
        \begin{block}{\textbf{Precisão} vs \textbf{Recall}}
            {\small Dois lados da mesma moeda:}

            \vspace{0.2cm}

            \textbf{Precisão:} ``Quando a IA diz que tem corrosão, ela está certa?''\\
            {\scriptsize $\rightarrow$ \textbf{85,6\%} das vezes sim}

            \vspace{0.3cm}

            \textbf{Recall:} ``A IA encontra toda a corrosão que existe?''\\
            {\scriptsize $\rightarrow$ Encontra \textbf{79,9\%} de toda corrosão}

            \vspace{0.3cm}

            \textbf{F1-Score:} média equilibrada dos dois\\
            {\scriptsize $\rightarrow$ \textbf{0,816} --- desempenho robusto}
        \end{block}
    \end{columns}
\end{frame}

\begin{frame}{Matriz de Confusão --- O ``Boletim'' do Modelo}
    \begin{block}{O que a IA acerta e o que ela erra}
        Leia como uma tabela de ``previsão vs realidade'':
    \end{block}

    \vspace{0.3cm}

    \begin{center}
    \begin{tikzpicture}[scale=0.85]
        % Headers
        \node[font=\small\bfseries] at (3.5,3.2) {Classe Real (Gabarito)};
        \node[font=\small\bfseries, rotate=90] at (-1.5,1) {Predição da IA};

        % Column headers
        \node[font=\scriptsize] at (1.5,2.6) {Nenhuma};
        \node[font=\scriptsize] at (3,2.6) {Leve};
        \node[font=\scriptsize] at (4.5,2.6) {Moderada};
        \node[font=\scriptsize] at (6,2.6) {Severa};

        % Row headers
        \node[font=\scriptsize, right] at (-1.2,2) {Nenhuma};
        \node[font=\scriptsize, right] at (-1.2,1.2) {Leve};
        \node[font=\scriptsize, right] at (-1.2,0.4) {Moderada};
        \node[font=\scriptsize, right] at (-1.2,-0.4) {Severa};

        % Diagonal (correct)
        \fill[verdepositivo!25] (0.9,1.6) rectangle (2.1,2.4);
        \fill[verdepositivo!25] (2.4,0.8) rectangle (3.6,1.6);
        \fill[verdepositivo!25] (3.9,0) rectangle (5.1,0.8);
        \fill[verdepositivo!25] (5.4,-0.8) rectangle (6.6,0);

        % Values
        \node[font=\small\bfseries, verdepositivo] at (1.5,2) {96\%};
        \node[font=\small\bfseries, verdepositivo] at (3,1.2) {93\%};
        \node[font=\small\bfseries, verdepositivo] at (4.5,0.4) {91\%};
        \node[font=\small\bfseries, verdepositivo] at (6,-0.4) {94\%};

        % Off-diagonal (errors)
        \node[font=\tiny, gray] at (3,2) {2\%};
        \node[font=\tiny, gray] at (4.5,2) {1\%};
        \node[font=\tiny, gray] at (6,2) {1\%};
        \node[font=\tiny, gray] at (1.5,1.2) {3\%};
        \node[font=\tiny, gray] at (4.5,1.2) {3\%};
        \node[font=\tiny, gray] at (6,1.2) {1\%};
        \node[font=\tiny, gray] at (1.5,0.4) {1\%};
        \node[font=\tiny, gray] at (3,0.4) {4\%};
        \node[font=\tiny, gray] at (6,0.4) {4\%};
        \node[font=\tiny, gray] at (1.5,-0.4) {0\%};
        \node[font=\tiny, gray] at (3,-0.4) {1\%};
        \node[font=\tiny, gray] at (4.5,-0.4) {5\%};

        % Grid
        \draw[gray!50] (0.9,2.4) grid[step=1.5] (6.6,-0.8);
        \draw[thick] (0.9,2.4) rectangle (6.6,-0.8);
    \end{tikzpicture}
    \end{center}

    \vspace{0.2cm}

    {\small \textcolor{verdepositivo}{\textbf{Diagonal verde}} = acertos. A IA erra mais entre classes \textit{vizinhas} (leve$\leftrightarrow$moderada), o que é esperado --- assim como um engenheiro hesitaria nos casos limítrofes.}
\end{frame}

% =============================================
% CONCLUSÃO
% =============================================
\section{Conclusões}

\begin{frame}{O Que Aprendemos}
    \begin{block}{Contribuições Deste Trabalho}
        \begin{enumerate}
            \item \textbf{Viabilidade comprovada:} IA consegue detectar e classificar corrosão em aço ASTM A572 Grau 50 com acurácia superior a 90\%

            \item \textbf{Attention U-Net superior:} o mecanismo de atenção melhora a segmentação em todas as métricas, especialmente em corrosão severa (+6,6\% em recall)

            \item \textbf{Transfer learning eficaz:} ResNet50 pré-treinada alcança 94,2\% de acurácia na classificação de severidade

            \item \textbf{Sistema dual:} a combinação classificação + segmentação oferece tanto triagem rápida quanto análise detalhada

            \item \textbf{Implementação 100\% em MATLAB:} acessível para engenheiros que já usam o ambiente
        \end{enumerate}
    \end{block}
\end{frame}

\begin{frame}{Limitações e Próximos Passos}
    \begin{columns}[T]
        \column{0.48\textwidth}
        \begin{alertblock}{Limitações Atuais}
            \begin{itemize}
                \item Dataset de 800 imagens (ampliável)
                \item Validação em condições controladas
                \item Específico para aço ASTM A572 Grau 50
                \item Necessita GPU para treinamento
            \end{itemize}
        \end{alertblock}

        \column{0.48\textwidth}
        \begin{exampleblock}{Trabalhos Futuros}
            \begin{itemize}
                \item Expandir para outros tipos de aço
                \item Integração com drones autônomos
                \item Aplicativo mobile para campo
                \item Monitoramento temporal (evolução)
                \item Validação em estruturas reais
            \end{itemize}
        \end{exampleblock}
    \end{columns}

    \vspace{0.5cm}

    \begin{block}{Mensagem Final}
        A Inteligência Artificial não substitui o engenheiro --- ela o \textbf{empodera} com ferramentas que tornam a inspeção mais segura, mais rápida, mais objetiva e mais acessível. O futuro da engenharia de manutenção estrutural é \textbf{assistido por IA}.
    \end{block}
\end{frame}

\begin{frame}{Publicações e Artigos Científicos}
    \begin{block}{Artigos Produzidos}
        \begin{enumerate}
            \item \textbf{Segmentação:} ``Comparative Analysis of U-Net and Attention U-Net for Automated Corrosion Detection in ASTM A572 Grade 50 Steel W-Beams'' --- Formato ASCE

            \item \textbf{Classificação:} ``Automated Corrosion Severity Classification in ASTM A572 Grade 50 Steel Using Deep Learning: A Hierarchical Approach for Structural Health Monitoring'' --- Formato ASCE/Elsevier
        \end{enumerate}
    \end{block}

    \vspace{0.3cm}

    \begin{block}{Referências Fundamentais}
        \begin{itemize}
            \item Ronneberger et al. (2015) --- U-Net
            \item Oktay et al. (2018) --- Attention U-Net
            \item He et al. (2015) --- ResNet
            \item Tan \& Le (2019) --- EfficientNet
            \item LeCun et al. (2015) --- Deep Learning (Nature)
        \end{itemize}
    \end{block}
\end{frame}

% =============================================
% SLIDE FINAL
% =============================================
\begin{frame}[plain]
    \begin{center}
        \vspace{1cm}
        {\Huge \textcolor{uerjazul}{\textbf{Obrigado!}}}

        \vspace{1cm}

        {\Large Perguntas?}

        \vspace{1cm}

        \begin{tabular}{rl}
            \faEnvelope & heitor.42540005@ucp.br \\
            \faUniversity & Universidade Católica de Petrópolis \\
            \faGithub & Código-fonte disponível mediante solicitação \\
        \end{tabular}

        \vspace{1cm}

        {\small Apresentação realizada na UERJ --- Faculdade de Engenharia}
    \end{center}
\end{frame}

\end{document}
